\section{Conclusion}\label{sec:conclusion}

We used simulations to answer one question: under realistic noise, how good can a Kerr-microcomb reservoir be, what limits it, and how
should it be compared fairly with software.

First, the encoding decides whether the limit can be approached. Giving every line of an input comb a random spectral phase before
modulating it lets the resonator operate in a coherently seeded regime, where the output is set by the input rather than by start-up
noise. Under the same physical conditions and selection budget this spectral-phase mask beats the single-pump encodings of the
literature and a uniformly modulated comb by one to two orders of magnitude on Mackey--Glass prediction, leads slightly on channel
equalisation and does not work on NARMA-10 with a single ring. What matters is the randomness of the line phases, so one modulator and
two spectral phase shapers suffice.

Second, the limit is set by noise, half of it at the source. With ideal detectors about half of the error floor comes from detuning
jitter and half from shot noise; another platform or more power does not lower it, and the readout-side changes (stronger drop
coupling, slower symbols) lower it in part. The jitter cannot be removed in the readout or avoided by another operating point, so the
intrinsic laser linewidth is a hard requirement: below a few kilohertz for our device, while an ordinary $100~\mathrm{kHz}$ laser makes
the error 5 to 6 times worse. The resonator contributes nonlinear features; memory has to come from external delay lines and readout
taps.

Third, in software terms the limit is below software of the same size. The best single ring matches a 31-node ESN fed through a
converter with about 7 effective bits, while at the resonator's rate of about $1~\mathrm{GBd}$ off-the-shelf converters offer about 9;
with the quantization of the resonator's 31 readings included, an ESN quantized only at its input is better at every depth from 6 to 12
bits. Reading several rings together lowers the error another 19 times, but large software reservoirs with lower input noise catch up
as well. The comparison is sensitive to the input noise given to software, so we report whole curves and suggest that future work do
the same.

The value of a Kerr-microcomb reservoir therefore does not lie in computing accuracy and can only come from the physical layer: a signal
that is optical to begin with, digitisation to be avoided, or optical latency and energy. This paper does not demonstrate it; doing so
needs a concrete optical application, an experiment, and a readout weighted in the optical domain that uses a single converter.
